% ------------------------------------------------------------
% English working translation layer.
% Source file: letters/Letter_Vasiu_2011.tex
% ------------------------------------------------------------

\documentclass[11pt,a4paper]{article}
\usepackage[utf8]{inputenc}
\usepackage[T1]{fontenc}
\usepackage{amsmath,amssymb,amsthm}
\usepackage{amsfonts}
\usepackage{mathtools}
\usepackage{geometry}
\usepackage{hyperref}
\usepackage{enumitem}
\usepackage{booktabs}
\usepackage{array}
\usepackage{mathrsfs}
\usepackage{xcolor}
\usepackage{microtype}
\usepackage{tikz-cd}

\geometry{a4paper, top=2.5cm, bottom=2.5cm, left=2.5cm, right=2.5cm}

\theoremstyle{plain}
\newtheorem{theorem}{Th\'eor\`eme}
\newtheorem{lemma}[theorem]{Lemma}
\newtheorem{proposition}[theorem]{Proposition}
\newtheorem{corollary}[theorem]{Corollary}
\newtheorem{conjecture}[theorem]{Conjecture}

\theoremstyle{definition}
\newtheorem{definition}[theorem]{D\'efinition}
\newtheorem{example}[theorem]{Example}
\newtheorem{remark}[theorem]{Remark}

\DeclareMathOperator{\Ext}{Ext}
\DeclareMathOperator{\Hom}{Hom}
\DeclareMathOperator{\Ker}{Ker}
\DeclareMathOperator{\Coker}{Coker}
\DeclareMathOperator{\Spec}{Spec}
\DeclareMathOperator{\Proj}{Proj}
\DeclareMathOperator{\Gal}{Gal}
\DeclareMathOperator{\Aut}{Aut}
\DeclareMathOperator{\End}{End}
\DeclareMathOperator{\Rep}{Rep}

\begin{document}
\begin{center}\small\emph{Literal English translation working draft. Formulae, numbering, and TeX structure are preserved; this is not yet a modern recast.}\end{center}
\vspace{0.5em}


\begin{flushright}
11/30/2011
\end{flushright}

\medskip

Dear Vasiu,

\medskip

Let $\mathcal{T}$ be a Tannakian category over an algebraically closed field $k$.
Your question amounts to: does $\mathcal{T}$ always have a fiber functor over $k$?
I believe it is true. The following argument is presumably too pedestrian,
hiding which ``compactness arguments'' are really relevant.

Let $A$ be the set of strictly full subcategories of $\mathcal{T}$, stable by
$\otimes$, subquotients and duals. Let $A_f \subset A$ be the set of those
generated by finitely many objects, using the same operations. For $\alpha \in A$,
I note $\mathcal{T}_\alpha$ the corresponding subcategory (as defined:
$\mathcal{T}_\alpha = \alpha$). The set $A$ is ordered by inclusion. I assume
we already know the existence and uniqueness, up to isomorphism, of fiber
functors for the $\mathcal{T}_\alpha$, $\alpha \in A_f$.

If $\mathcal{T}_\alpha \subset \mathcal{T}_\beta$, it makes sense to say that
a fiber functor $\omega_\beta$ one $\mathcal{T}_\beta$ extends (or the nose) a
fiber functor $\omega_\alpha$ one $\mathcal{T}_\alpha$. If an extension, up to
isomorphism, exists, an actual extension exists too.

Let us order the set of $(\alpha, \omega_\alpha)$: $\alpha \in A$,
$\omega_\alpha$ fiber functor one $\mathcal{T}_\alpha$, by ``$\omega_\beta$
extends $\omega_\alpha$''. To avoid set theoretical difficulties, one could
consider only the fiber functors $\omega_\alpha$ such that the
$\omega_\alpha(X)$ take values in the set of vector spaces $k^n$, $n \in \mathbb{N}$.

The ordered set of $(\alpha, \omega_\alpha)$ is inductive (Bourbaki Ens Ch~3
\S2.4): if $I$ is a totally ordered subset, the ``union'' of the
$(\alpha, \omega_\alpha)$ is a $\mathcal{T}_0$ in $A$
$(= \bigcup_{(\alpha,\omega_\alpha) \in I} \mathcal{T}_\alpha)$ with the fiber
functor $\omega_0$ characterized by $\omega_0|\mathcal{T}_\alpha = \omega_\alpha$.
This ``union'' majorizes $I$. By Bourbaki loc.~cit.~Th~2, the ordered set of the
$(\alpha, \omega_\alpha)$ has a maximal element: $(\mathcal{T}_1, \omega_1)$.
To prove that $\mathcal{T}_1 = \mathcal{T}$, it suffices to prove the following

\medskip
\noindent\textbf{Lemma 1.} Let $\mathcal{T}'$ be in $A$ and $\mathcal{T}''$ be in
$A_f$. Let $\langle \mathcal{T}', \mathcal{T}'' \rangle$ be generated by
$\mathcal{T}'$ and $\mathcal{T}''$; it is in $A$. Then, any fiber functor
$\omega$ one $\mathcal{T}'$ can be extended (one the nose, or up to isomorphism,
this amounts to the same) to $\langle \mathcal{T}', \mathcal{T}'' \rangle$.

\medskip
\noindent I first prove

\medskip
\noindent\textbf{Lemma 2.} Assumed $\mathcal{T}'$ is in $A_f$ too.
``Restriction'' is then an equivalence of categories
\[
\text{(fiber functors on } \langle \mathcal{T}', \mathcal{T}'' \rangle \text{)}
\longrightarrow
\text{(triples of)}
\]
a fiber functor $\omega'$ one $\mathcal{T}'$, $\omega''$ one $\mathcal{T}''$ and an
isomorphism $\tau$ of the restrictions of $\omega'$ and $\omega''$ to
$\mathcal{T}' \cap \mathcal{T}''$ (also in $A_f$).

\medskip

We may assume that $\langle \mathcal{T}', \mathcal{T}'' \rangle$ is the category
of representations of a group $G$, and that for invariant subgroups $A$ and $B$,
$\mathcal{T}'$ (resp.~$\mathcal{T}''$) is the subcategory of representations
where $A$ (resp.~$B$) acts trivially. That they generate
$\langle \mathcal{T}', \mathcal{T}'' \rangle$ means that
$A \cap B = \{e\}$. The intersection $\mathcal{T}' \cap \mathcal{T}''$ is the
category of representations one which $AB$ acts trivially.

The triples $(\omega', \omega'', \tau)$ are all isomorphic: as all $\omega'$
(resp.~all $\omega''$) are isomorphic, it suffices to see that
$(\omega', \omega'', \tau_1)$ and $(\omega', \omega'', \tau_2)$ are isomorphic.
Indeed $\tau_1$ and $\tau_2$ differ by an automorphism of
$\omega'|\mathcal{T}' \cap \mathcal{T}''$, and such an automorphism lifts to
an automorphism of $\omega'$:
\[
G/A(k) \longrightarrow G/AB(k) \quad \text{is onto.}
\]

We hence have here categories with just one isomorphism class of objects, and
the question is to compare automorphism groups. We need to check
\[
G \;\xrightarrow{\sim}\; \{(g', g'') \in G/A, G/B \mid g' \text{ and } g''
\text{ have same image in } G/AB\}
\]
\[
@@MATH0@@
\]
Injectivity of this morphism of groups amounts to $A \cap B = \{e\}$.

Surjectivity: if $\tilde{g}', \tilde{g}''$ lift $g', g''$,
$\tilde{g}'' = \tilde{g}' \cdot a \cdot b$ and
$\tilde{g}'' \cdot b = \tilde{g}' \cdot a \in G$ maps to $(g', g'')$.

\medskip
\noindent\textbf{Lemma 3.} Same as lemma 2, but $\mathcal{T}'$ only assumed to
be in $A$.

\medskip

Let $B$ be the set of $\mathcal{T}_\alpha$ ($\alpha \in A_f$) contained in
$\mathcal{T}'$. one has
$\langle \mathcal{T}', \mathcal{T}'' \rangle =
\bigcup_{\beta \in B} \langle \mathcal{T}_\beta, \mathcal{T}'' \rangle$.
one has equivalences
\[
\text{(fiber functors on } \mathcal{T} \text{)}
\;\xrightarrow{\sim}\;
\text{(fiber functors on the } \mathcal{T}_\beta\text{,}
\]
more a compatible system of isomorphisms
$\omega_\beta|\mathcal{T}_\gamma \xrightarrow{\sim} \omega_\gamma$ for
$\mathcal{T}_\gamma \subset \mathcal{T}_\beta$; compatible: condition for
$\mathcal{T}_\delta \subset \mathcal{T}_\gamma \subset \mathcal{T}_\beta$).

Same for fiber functors one the
$\langle \mathcal{T}', \mathcal{T}'' \rangle =
\bigcup \langle \mathcal{T}_\beta, \mathcal{T}'' \rangle$.

The $\mathcal{T}_\beta \cap \mathcal{T}''$ have a largest element: they if
$\mathcal{T}'' = \Rep(G'')$, they correspond to invariant subgroups of $G''$,
subgroups are closed subschemes, and one uses the noetherian property. If
$\beta_0$ is such that $\mathcal{T}_{\beta_0} \cap \mathcal{T}''$ is the
largest $\mathcal{T}_\beta \cap \mathcal{T}''$, for any $\beta > \beta_0$,
extending $\omega_\beta = \omega|\mathcal{T}_\beta$ to
$\langle \mathcal{T}_\beta, \mathcal{T}'' \rangle$ amounts to extending
$\omega_\beta|\mathcal{T}_\beta \cap \mathcal{T}''$ to $\mathcal{T}''$
(lemma~2), that is to extend $\omega_{\beta_0}$ from
$\mathcal{T}_{\beta_0} \cap \mathcal{T}''$ to $\mathcal{T}''$. If we choose
one such extension, we get up to unique isomorphism a system of extensions of
the $\omega_\beta$ to $\langle \mathcal{T}_\beta, \mathcal{T}'' \rangle$, and
by gluing them an extension of $\omega$ to
$\langle \mathcal{T}', \mathcal{T}'' \rangle$.

This is all we need to conclude that $\mathcal{T}' = \mathcal{T}$.

This proof should be cleaned up. After all, we are proving that some projective
system of gerbs (of the fiber functors one the $\mathcal{T}_\alpha$,
$\alpha \in A_f$), where ``projective system'' is taken in a 2-categorical
sense, has a not empty projective limit (again limit in a 2-categorical sense).
Maybe such a translation would make the ``compactness arguments'' used cleaner.
They were two of them: ``fiber functor is a property of finite type (cf
Bourbaki Ens Ch~3 \S4.5)'', and the noetherian property for subgroups.

The same argument gives uniqueness up to isomorphism of $\omega$, over $k$; we
get a maximal $\mathcal{T}_\alpha$ ($\alpha \in A$) over which we have an
isomorphism, and extend further if $\mathcal{T}_\alpha \neq \mathcal{T}$.

\medskip
\medskip

\hfill Best,

\hfill P.~Deligne

\end{document}
